\section{Sensitivity and estimated difference means}\label{sec:plugin}

\begin{theorem}[Smoothness sensitivity and conditioning]\label{thm:sensitivity}
For each interior target $s\in(0,1)$, let
$b_j(s)=w_j(s)[1-w_j(s)]$ and
$B(s)=r^{-1}\sum_{j=1}^r b_j(s)>0$.
Then the spectral weights and the fixed-$\mu$ fitted trend
are differentiable and satisfy
\begin{align}
\frac{\mathrm d w_j}{\mathrm ds}
&=\frac{b_j(s)}{B(s)},\quad
0<\frac{\mathrm dw_j}{\mathrm ds}\le r,
\qquad \sum_{j=1}^r\frac{\mathrm dw_j}{\mathrm ds}=r,
\label{eq:wdot}\\
\frac{\partial\widehat\tau_{s,\mu}}{\partial s}
&=-V\diag\!\left(\frac{b_1(s)}{B(s)},\ldots,
\frac{b_r(s)}{B(s)}\right)c_\mu.\label{eq:trenddot}
\end{align}
For all $s,t\in[0,1]$,
\begin{equation}\label{eq:lipschitz}
\norm{\widehat\tau_{s,\mu}-\widehat\tau_{t,\mu}}_2
\le r|s-t|\norm{c_\mu}_2 .
\end{equation}
For fixed $s$ and two numeric means $\mu_1,\mu_2$,
\begin{equation}\label{eq:mu-lipschitz}
\norm{\widehat\tau_{s,\mu_1}-\widehat\tau_{s,\mu_2}}_2
=|\mu_1-\mu_2|
\left[\sum_{j=1}^r
\frac{w_j(s)^2(U^\top\one_r)_j^2}{\sigma_j^2}\right]^{1/2}.
\end{equation}
Moreover the derivative of the scalar spectral map obeys
\begin{equation}\label{eq:variance-id}
G'(x_s)=s(1-s)
-\frac1r\sum_{j=1}^r[w_j(s)-s]^2
\le s(1-s),
\end{equation}
and hence $\mathrm dx_s/\mathrm ds\ge[s(1-s)]^{-1}$.
\end{theorem}
\begin{proof}
Differentiate $G(x_s)=s$ and use
$G'(x_s)=B(s)>0$ to obtain
$\mathrm dx_s/\mathrm ds=1/B(s)$.
The chain rule then gives
$\mathrm dw_j/\mathrm ds=b_j(s)/B(s)$.
Because $rB(s)=\sum_k b_k(s)\ge b_j(s)$,
the derivative is at most $r$; summing the derivatives gives $r$.
Differentiate \eqref{eq:direct-fit} with $y,\mu$ fixed for
\eqref{eq:trenddot}. The mean-value theorem and
\eqref{eq:wdot}, extended to endpoint limits by continuity,
yield $|w_j(s)-w_j(t)|\le r|s-t|$.
Since $V^\top V=I_r$,
\[
\norm{V\diag(w(s)-w(t))c_\mu}_2
\le\max_j|w_j(s)-w_j(t)|\norm{c_\mu}_2,
\]
which proves \eqref{eq:lipschitz}.
The $\mu$-dependence of \eqref{eq:direct-fit} is affine:
\[
\widehat\tau_{s,\mu_1}-\widehat\tau_{s,\mu_2}
=(\mu_1-\mu_2)V\diag(w(s))\Sigma^{-1}U^\top\one_r.
\]
Its Euclidean norm equals the right side of
\eqref{eq:mu-lipschitz} because $V$ is an isometry on $\R^r$.
Finally, using $s=r^{-1}\sum_jw_j$,
\[
G'(x_s)=\frac1r\sum_j(w_j-w_j^2)
=s-\left(s^2+\frac1r\sum_j(w_j-s)^2\right),
\]
which gives \eqref{eq:variance-id}. Inverting the strictly
positive derivative establishes the lower bound on
$\mathrm dx_s/\mathrm ds$.
\end{proof}

\begin{remark}[Conditioning is not eliminated]\label{rem:cond}
The factor $\mathrm dx_s/\mathrm ds$ becomes large toward both
endpoints. This is also the conditioning of a log-penalty inverse,
up to the constant shift in \eqref{eq:implicit}. In contrast,
\eqref{eq:wdot} bounds the sensitivity of each \emph{weight} to
the requested smoothness. A well-conditioned weight curve does not
imply that its scalar log-odds root is well-conditioned, or that
an inaccurate computed SVD can be repaired by reparameterization.
\end{remark}

\begin{proposition}[Endpoint penalty scales]\label{prop:asymptotic}
With $n,d$ and the positive spectrum fixed, define
$A=r^{-1}\sum_{j=1}^r\delta_j>0$ and
$B=r^{-1}\sum_{j=1}^r\delta_j^{-1}>0$.
Then
\begin{align}
S_d(\lambda)&=A\lambda+O(\lambda^2)
 &&(\lambda\downarrow0),\\
1-S_d(\lambda)&=B/\lambda+O(\lambda^{-2})
 &&(\lambda\uparrow\infty).
\end{align}
Consequently $\Lambda_d(s)=s/A+O(s^2)$ as $s\downarrow0$
and $\Lambda_d(s)=B/(1-s)+O(1)$ as $s\uparrow1$.
\end{proposition}
\begin{proof}
For each fixed $\delta_j>0$,
$\lambda\delta_j/(1+\lambda\delta_j)
=\lambda\delta_j+O(\lambda^2)$ as $\lambda\to0$.
The finite average gives the first expansion. Similarly
$1/(1+\lambda\delta_j)
=(\lambda\delta_j)^{-1}+O(\lambda^{-2})$
as $\lambda\to\infty$, yielding the second.
Strict invertibility from Theorem~\ref{thm:smoothness} allows
termwise inversion of each first-order expansion. Explicitly,
$s=A\lambda+O(\lambda^2)$ implies
$\lambda=s/A+O(s^2)$; and with $t=1-s$,
$t=B/\lambda+O(\lambda^{-2})$ implies
$1/\lambda=t/B+O(t^2)$ and hence
$\lambda=B/t+O(1)$.
\end{proof}

\subsection{Observed-difference plug-in mean}
The previous results treat $\mu$ as fixed. The option
\begin{equation}\label{eq:plug}
\widehat\mu(y)=\frac1r\one_r^\top D_dy
=:g^\top y,\qquad g=D_d^\top\one_r/r,
\end{equation}
instead defines a different \emph{data-dependent linear mapping}
$y\mapsto\widehat\tau_{s,\widehat\mu(y)}$ for every fixed target $s$.
Its full Jacobian need not equal $H_s$.

\begin{theorem}[Exact plug-in smoother and EDF correction]\label{thm:plugin}
Let $s\in[0,1]$ be fixed and put
$a=U^\top\one_r$,
$b_s=V\diag(w(s))\Sigma^{-1}a$.
The estimator obtained by substituting \eqref{eq:plug} into
\eqref{eq:direct-fit} is the linear map
\begin{equation}\label{eq:plug-hat}
\widehat\tau_{s,\widehat\mu(y)}
=T_sy,\qquad
T_s=H_s+b_sg^\top,
\end{equation}
where $H_s=I_n-V\diag(w(s))V^\top$.
Its \emph{unconditional} trace (holding $s$ fixed, but
not holding the estimated mean fixed) is
\begin{equation}\label{eq:edf-plug}
\begin{aligned}
\nu_{\rm plug}(s)&=\tr T_s
=n-rs+\Delta(s),\\
\Delta(s)&=\frac1r\sum_{j=1}^r w_j(s)a_j^2,
\qquad 0\le\Delta(s)\le1.
\end{aligned}
\end{equation}
In particular,
\[
\nu_{\rm plug}(0)=n,\qquad
\nu_{\rm plug}(1)=d+1,
\]
whereas the conditional fixed-$\mu$ trace equals $n-rs$,
with limiting value $d$ at $s=1$.
\end{theorem}
\begin{proof}
Equation \eqref{eq:direct-fit} can be written
$\widehat\tau_{s,\mu}=H_sy+\mu b_s$.
Replacing $\mu$ by $g^\top y$ gives \eqref{eq:plug-hat}.
Using $\tr(b_sg^\top)=g^\top b_s$ and
$D_d=U\Sigma V^\top$,
\[
\begin{aligned}
g^\top b_s
&=\frac1r\one_r^\top D_dV\diag(w)\Sigma^{-1}a\\
&=\frac1r\one_r^\top U\Sigma\diag(w)\Sigma^{-1}U^\top\one_r\\
&=\frac1r\,a^\top\diag(w)a
=\frac1r\sum_jw_ja_j^2.
\end{aligned}
\]
Together with $\tr H_s=n-rs$ this proves
\eqref{eq:edf-plug}. Since $U$ is orthogonal,
$\sum_j a_j^2=\norm{\one_r}_2^2=r$.
Every $w_j\in[0,1]$, so $\Delta\in[0,1]$.
At $s=0$ all weights are zero and
$\Delta(0)=0$; at $s=1$ every weight is one,
giving $\Delta(1)=1$. The endpoint
values follow.
\end{proof}

\begin{corollary}[A caution for interpreting normalized smoothness]\label{cor:edf}
If $\mu$ is estimated from observed differences, the index $s$
in \eqref{eq:S} still specifies the \emph{conditional} trace
reduction $r^{-1}\sum_jw_j(s)$. The reduction calculated from
the full plug-in operator is instead
\[
\frac{n-\tr(T_s)}r=s-\frac{\Delta(s)}r.
\]
These two quantities differ for interior $s$, and their endpoint
values at $s=1$ are $1$ and $1-1/r$, respectively.
\end{corollary}
\begin{proof}
Substitute \eqref{eq:edf-plug}. For $0<s<1$, the
$w_j(s)$ are strictly positive, while $a\ne0$ because
$U$ is orthogonal and $\one_r\ne0$, so $\Delta(s)>0$.
The endpoint values follow from Theorem~\ref{thm:plugin}.
\end{proof}

\begin{corollary}[Degenerate one-difference case]\label{cor:plugin-r1}
If $r=n-d=1$ and $\widehat\mu(y)=D_dy$, then
$\widehat\tau_{s,\widehat\mu(y)}=y$ for \emph{every}
$s\in[0,1]$. In particular, the full plug-in trace
is $n$, even though the conditional trace equals $n-s$.
\end{corollary}
\begin{proof}
For $\tau=y$ the residual in \eqref{eq:objective} is zero,
and the penalty $D_dy-\widehat\mu(y)$ also vanishes.
Thus $J_{\lambda,\widehat\mu(y)}(y)=0$ for every finite
$\lambda\ge0$, the global minimum of the nonnegative
objective. Uniqueness from Theorem~\ref{thm:pls}
gives $\widehat\tau=y$. The limit $s=1$ is also $y$
by continuity; differentiating the identity map gives trace $n$.
Alternatively, since there is only one positive spectral
weight, Theorem~\ref{thm:plugin} gives
$\Delta(s)=w_1(s)=s$ and
$\nu_{\rm plug}(s)=n-s+s=n$.
\end{proof}

\begin{remark}[Degrees of freedom versus statistical claims]
The trace $\tr(T_s)$ is the exact divergence of the
linear fitted-value map for a fixed $s$.
Interpreting this divergence as a covariance-based
effective number of parameters requires a statistical model
(e.g., homoskedastic independent Gaussian errors).
If $s$ itself is selected using $y$, further
data-dependence arises, and neither trace formula
automatically describes the full adaptive procedure.
Similarly, deleting one observation and holding a
plug-in $\widehat\mu$ fixed is not an exact leave-one-out
refit of the complete plug-in estimator.
\end{remark}
